\section{Stage II: whole-slide neural aggregation}
\begin{figure}[!htbp]
\centering
\resizebox{\columnwidth}{!}{\begin{tikzpicture}[
  node distance=0.45cm and 0.6cm,
  box/.style={draw, rounded corners=2pt, align=center, font=\scriptsize, minimum width=2.4cm, minimum height=0.6cm}
]
\node[box, fill=green!8] (bags) {Patch feature bag};
\node[box, fill=green!12, below left=of bags] (a) {Branch A\\TransMIL-style global correlation};
\node[box, fill=green!12, below right=of bags] (b) {Branch B\\nnMIL-style gated attention};
\node[box, fill=blue!10, below=of bags, yshift=-1.0cm] (tokens) {Branch tokens\\$[\,b_A \,;\, b_B\,]$};
\node[box, fill=blue!14, below=of tokens] (fusion) {Corrected branch-token\\self-attention fusion};
\node[box, fill=green!8, below=of fusion] (pool) {Hierarchical pooling\\topology branch $\cdot$ pruning};
\node[box, below=of pool] (head) {Slide-level head (QWK)};
\draw[->] (bags) -- (a);
\draw[->] (bags) -- (b);
\draw[->] (a) -- (tokens);
\draw[->] (b) -- (tokens);
\draw[->] (tokens) -- (fusion);
\draw[->] (fusion) -- (pool);
\draw[->] (pool) -- (head);
\end{tikzpicture}}
\caption{TransnnMIL architecture family. The repaired canonical model fuses
explicit branch tokens by self-attention. Historical fusion and topology
results are withdrawn; post-repair PANDA runs are preliminary and pending
matched controlled reruns (see Section~5).}
\label{fig:transnnmil}
\end{figure}


The second stage of PA-NF converts patch-level representations into slide-level predictions. A whole-slide image can contain tens of thousands of candidate tissue patches, so the learning problem is not only feature quality but how information is selected, combined, and propagated across the bag. TransnnMIL is the custom architecture family developed for this stage.

\subsection{TransnnMIL architecture}
The canonical implementation combines two complementary aggregation paths. A TransMIL-style branch models global correlations among patch embeddings \cite{shao2021transmil}, while an nnMIL/attention-style branch learns gated instance aggregation \cite{ilse2018attention}. Their slide representations are projected into explicit branch tokens and fused with self-attention over the branch-token sequence. This corrected design allows genuine interaction between the branches rather than reducing to a single-position attention operation.

The wider TransnnMIL family also contains:
\begin{itemize}
\item hierarchical spatial pooling over learned regions;
\item an optional topology branch with GATv2, GraphSAGE, or GIN processing;
\item coordinate-aware k-nearest-neighbour graph construction and HDF5 graph caching;
\item adaptive patch-pruning experiments;
\item concat, gate, and learned branch-attention fusion controls; and
\item two-branch and three-branch variants for matched ablation studies.
\end{itemize}

These components are evaluated under frozen, architecture-matched controls rather than inferred to be beneficial from their design alone.

\subsection{PANDA feature substrate}
PANDA \cite{bulten2022panda} provides the primary slide-level test bed. Of 10,614 candidate feature bags, 10,611 were readable. They contain 768-dimensional Phikon patch representations; the inherited full-feature split comprised 8,488 training and 2,123 validation slides. The repaired-model comparison used a separately frozen train/selection/confirmation protocol. Earlier pre-remediation optimization records are retained in the supplement rather than used as evidence for this repaired comparison.

\subsection{Completed repaired-architecture matched PANDA comparison}
The subsequent frozen experiment compared seven architectures at five matched seeds (7, 19, 42, 123, 2025), for 35 completed fits on the PANDA Phikon feature bags. It used a slide-disjoint 70/15/15 train/selection/confirmation partition, with checkpoint choice on selection rather than confirmation data. Patient-level grouping has not been established.

\begin{table}[t]
\centering
\caption{Completed repaired-architecture comparison: five-seed mean confirmation quadratic weighted kappa (QWK). This is distinct from the pre-remediation optimizer study.}
\label{tab:transnnmil_repaired_matched}
\small
\begin{tabular}{lr}
\toprule
Frozen architecture & Mean confirmation QWK \\
\midrule
TransMIL & 0.805826 \\
TransnnMIL gate & 0.798476 \\
TransnnMIL learned branch attention & 0.795407 \\
TransnnMIL concat & 0.794834 \\
AttentionMIL & 0.794186 \\
Repaired TransnnMIL (primary) & 0.792195 \\
nnMIL & 0.779618 \\
\bottomrule
\end{tabular}
\end{table}

The primary repaired-minus-TransMIL QWK difference was $-0.013631$, with zero of five paired seeds favoring the repaired model; the registered primary comparison therefore failed. The repaired model exceeded nnMIL by $+0.012577$ on average (five of five seeds) but this does not reverse the failed primary. In particular, the full registered hierarchical bootstrap criterion is not established by five mean checkpoint scores alone. The gate variant included one seed flagged for practical branch collapse. These results demonstrate completed training and a negative primary comparator outcome, not superiority of repaired fusion.

\subsection{Spatial tissue-dynamics extension}
PA-NF also contains an experimental WSI-NCA line that replaces one-shot bag aggregation with repeated local state updates over spatial neighbourhoods. A frozen PANDA-300 ladder compared static, one-step, four-step tied, shuffled-topology, and untied controls across five seeds. The four-step tied model did not outperform the static or one-step controls on mean held-out QWK, and real topology did not consistently outperform shuffled topology. WSI-NCA therefore remains an active spatial-modeling extension rather than a demonstrated predictive improvement in the current pipeline.
